\newcommand{\dfonc}[1]{\mathscr{D}_{[#1]}}
%\newcommand{\operator}[1]{\hat{\bm{#1}}} % pour les operaeur vecteur
\newcommand{\operator}[1]{\hat{\boldsymbol{#1}}}

\definecolor{Prune}{RGB}{99,0,60} % l14-33 :

% --- Bois (palette principale)
\definecolor{woodNoyer}{HTML}{3A2A1F}
\definecolor{woodChene}{HTML}{4A3628}
\definecolor{woodAcajou}{HTML}{5B2C2C}

% --- Verts profonds
\definecolor{greenForet}{HTML}{0B3D2E}
\definecolor{greenSapin}{HTML}{0F4C3A}

% --- Neutres
\definecolor{grayPierre}{HTML}{4A4A4A}
\definecolor{grayBrume}{HTML}{7A7A7A}

% --- Custom colors
\definecolor{cvred}{HTML}{B71C1C}
\definecolor{headerblue}{HTML}{1565C0}
\definecolor{cvgreen}{HTML}{2E7D32}

\definecolor{titleblue}{HTML}{4a7aa4}

\definecolor{colorStitre}{RGB}{140,0,0}
\definecolor{colorSStitre}{RGB}{190,140,140}
\definecolor{colorSSStitre}{RGB}{140,140,140}
\def\colorStitre{colorStitre}
\def\colorSStitre{colorSStitre}
\def\colorSSStitre{colorSSStitre}

\definecolor{colorOne}{HTML}{443E46}
%\definecolor{colorTwo}{HTML}{F6DEB8}
\definecolor{colorTwo}{HTML}{FFFFFF}
\definecolor{colorThree}{HTML}{908CA4}
\definecolor{colorFour}{HTML}{57659E}
\definecolor{colorFive}{HTML}{C57284}
\definecolor{colorSix}{HTML}{FF5B69}
\definecolor{colorGold}{HTML}{FFD700}

% Raccourcis pour les couleurs
\def\colorOne{colorOne}
\def\colorTwo{colorTwo}
\def\colorThree{colorThree}
\def\colorFour{colorFour}
\def\colorFive{colorFive}
\def\colorSix{colorSix}
\def\colorGold{colorGold}

\makeglossaries

\newacronym{LL}{LL}{Lieb–Liniger}
\newacronym{NS}{NLS}{Nonlinear Schrödinger}
\newacronym{GP}{GP}{Gross–Pitaevskii}
\newacronym{GGE}{GGE}{Generalized Gibbs Ensemble}
\newacronym{GE}{GE}{Gibbs Ensemble}
\newacronym{BA}{BA}{Bethe Ansatz}
\newacronym{FBA}{FBA}{Fermionic Bethe Ansatz}
\newacronym{TBA}{TBA}{Thermodynamic Bethe Ansatz}
\newacronym{qBEC}{qBEC}{quasi Bose–Einstein Condensate}
\newacronym{GHD}{GHD}{Generalized Hydrodynamics}
\newacronym{LDA}{LDA}{Local Density Approximation}
\newacronym{DMD}{DMD}{Dispositif de Micromirroirs Digitaux}
\newacronym{PMO}{PMO}{Piège Magnéto-Optique}
\newacronym{DC}{DC}{Direct Current}
\newacronym{AC}{AC}{Alternating Current}
\newacronym{SYRTE}{SYRTE}{Systèmes de Référence Temps-Espace}
\newacronym{DFB}{DFB}{Distributed Feedback}
\newacronym{DBR}{DBR}{Distributed Bragg Reflector}
\newacronym{TA}{TA}{Tapered Amplifier}
\newacronym{AOM}{AOM}{Acousto-Optic Modulator}
\newacronym{PBS}{PBS}{Polarizing Beam Splitter}
\newacronym{LCF}{LCF}{Laboratoire Charles Fabry}
\newacronym{C2N}{C2N}{Centre de nanosciences et de nanotechnologies}
\newacronym{BCB}{BCB}{benzocyclobuténe}
\newacronym{SiC}{SiC}{silicon carbide}
\newacronym{Au}{Au}{aurum}
%\newacronym{AOM}{AOM}{modulateur acousto-optique}
\newacronym{CCD}{CCD}{Charge-Coupled Device}
\newacronym{TF}{TF}{Thomas-Fermi}
\newacronym{CESQ}{CESQ}{Centre européen de sciences quantiques}
\newacronym{QED}{QED}{Quantum Electrodynamics}
\newacronym{BCS}{BCS}{Bardeen–Cooper–Schrieffer}
\newacronym{YY}{YY}{Yang–Yang}
\newacronym{TG}{TG}{Tonks-Girardeau}
\newacronym{MFT}{MFT}{Macroscopic Fluctuation Theory}
\newacronym{GOE}{GOE}{Gaussian Orthogonal Ensemble}
\newacronym{GUE}{GUE}{Gaussian Unitary Ensemble}
\newacronym{RMT}{RMT}{Random Matrix Theory}
\newacronym{ABA}{ABA}{Algebraic Bethe Ansatz}

%\usepackage[french]{babel}
%\usetikzlibrary{babel}
\tikzset{
  every picture/.style={
      execute at begin picture={\shorthandoff{:;!?}}
    }
}

\tikzstyle{every picture}+=[remember picture]
\tikzstyle{na} = [shape=ellipse,inner sep=0pt]

% \newcommand{\sectiontikz}[1]{

%   \noindent
%   \begin{tikzpicture}[baseline=(title.base)]

%     \node[
%       inner sep=2pt
%     ] (title)
%     {
%       \fontsize{24}{24}\selectfont
%       \bfseries
%       #1
%     };

%     \draw[titre, line width=5pt]
%     ([xshift=-6cm]title.west) --
%     ([xshift=-0.7cm]title.west);

%     \draw[titre, line width=5pt]
%     ([xshift=0.7cm]title.east) --
%     ([xshift=6cm]title.east);

%   \end{tikzpicture}

% }

% \newcommand{\sectiontikz}[1]{
%   \vspace{-0.5cm}
%   \noindent
%   \begin{tikzpicture}

%     % Titre centré
%     \node[
%       inner sep=2pt,
%       %text width=0.7\linewidth,
%       %align=justify,
%       align=center, text width=min(0.7\linewidth,\widthof{
%         \fontsize{15}{10}\selectfont \bfseries {\Roman{section}}\ -\ #1}),
%       font=\fontsize{15}{10}\selectfont \bfseries, ] (title) at (\linewidth/2,0) {
%       {\Roman{section}}\ -\ #1 };%color = colorStitre 

%     % Ligne gauche
%     \draw[colorStitre, line width=5pt]
%     (0,0) --
%     ([xshift=-0.7cm]title.west);

%     % Ligne droite
%     \draw[colorStitre, line width=5pt]
%     ([xshift=0.7cm]title.east) --
%     (\linewidth,0);

%   \end{tikzpicture}

% }

\newcommand{\sectiontikz}[1]{
  %\vspace{-0.5cm}
  \noindent
  \begin{tikzpicture}

    \node[
      inner sep=0pt,
      %text width=0.7\linewidth,
      %align=justify,
      align=center, font=\fontsize{15}{10}\selectfont \bfseries, color = colorStitre
    ] (title) at (\linewidth/2,0) {
      \begin{varwidth}{0.7\linewidth}
        \centering
        {\Roman{section}}\ -\ #1
      \end{varwidth}
    };

    \draw[colorStitre,line width=0.3cm]
    (0,0)
    --
    ([xshift=-0.5cm]title.west);

    \draw[colorStitre,line width=0.3cm]
    ([xshift=0.5cm]title.east)
    --
    (\linewidth,0);

  \end{tikzpicture}

}

\titleformat{\section}
{\normalfont}
{}
{-1.8mm}
{\sectiontikz}
[\vspace{-1.0cm}]

%

%\titleformat{\section}{FORMAT}{LABEL}{SEP}{CODE}
%\titleformat{commande}[forme]{format}{label}{sep}{before-code}[after-code]\thesection

%\titlespacing*{\section}{0pt}{1cm}{0.7cm}

\newcommand{\subsectiontikz}[1]{
  %\vspace{-0.5cm}
  \noindent
  \begin{tikzpicture}

    \node[
      inner sep=0pt,
      %text width=0.7\linewidth,
      %align=justify,
      align=center, font=\fontsize{12}{10}\selectfont \bfseries, color = colorSStitre
    ] (title) at (\linewidth/2,0) {
      \begin{varwidth}{0.7\linewidth}
        \centering
        \thesubsection\ -\ #1
      \end{varwidth}
    };

    \draw[colorSStitre,line width=0.15cm]
    (0,0)
    --
    ([xshift=-0.3cm]title.west);

    \draw[colorSStitre,line width=0.15cm]
    ([xshift=0.3cm]title.east)
    --
    (\linewidth,0);

  \end{tikzpicture}

}

\titleformat{\subsection}
{\normalfont}
{}
{0pt}
{\subsectiontikz}
[\vspace{-1.0cm}]

\newcommand{\subsubsectiontikz}[1]{
  %\vspace{-0.5cm}
  \noindent
  \begin{tikzpicture}

    \node[
      inner sep=0pt,
      %text width=0.7\linewidth,
      %align=justify,
      align=center, font=\fontsize{10}{10}\selectfont \bfseries, color =
      colorSSStitre ] (title) at (\linewidth/2,0) {
      \begin{varwidth}{0.7\linewidth}
        \centering
        \thesubsubsection\ -\ #1
      \end{varwidth}
    };

    \draw[colorSSStitre,line width=0.10cm]
    (0,0)
    --
    ([xshift=-0.2cm]title.west);

    \draw[colorSSStitre,line width=0.10cm]
    ([xshift=0.2cm]title.east)
    --
    (\linewidth,0);

  \end{tikzpicture}

}

\titleformat{\subsubsection}
{\normalfont}
{}
{0pt}
{\subsubsectiontikz}
[\vspace{-1.0cm}]

%\newcommand{\ptFleche}[2]{		% Déclaration d'une extrémité de flèche
%    \shorthandoff{:;!?}% <- désactive les shorthands
%    \tikz[baseline=(#1.base)]\node[na](#1){#2};
%  }

\newcommand{\ptFleche}[2]{%
  \tikz[baseline=(#1.base)]{
    \shorthandoff{:;!?}
    \node[na](#1){#2};
  }%
}

\newcommand{\Fleche}[5][thick]{%
  \begin{tikzpicture}[overlay]
    \path[->,#1](#2) edge [out=#4, in=#5] (#3);
  \end{tikzpicture}
}

% --- Commande FlecheEllipse ---
% #1 = nœud source (nom)
% #2 = angle sortie du nœud source
% #3 = angle entrée du nœud cible
% #4 = décalage dx, dy sous forme (dx,dy) en cm
% #5 = couleur ellipses et flèche
% #6 = texte du nœud cible B
\newcommand{\FlecheEllipseo}[6]{%
  \begin{tikzpicture}[overlay, remember picture]
    % Coordonnées source
    \path let \p1 = (#1) in
    coordinate (Acoord) at (\x1,\y1);

    % Coordonnées cible
    \pgfmathsetmacro{\dx}{#4}
    \pgfmathsetmacro{\dy}{#5} % ici on sépare les deux pour le calcul

    \coordinate (B) at ([shift={(#4)}]#1);

    % Nœud A entouré d'ellipse
    \node[draw,#5,ellipse,inner sep=2pt] (#1-node) at (#1) {#1};
    % Nœud B entouré d'ellipse
    \node[draw,#5,ellipse,inner sep=2pt] (B-node) at (B) {#6};
    % Flèche de A vers B
    \draw[->,#5,thick] (#1-node) edge [out=#2,in=#3] (B-node);
  \end{tikzpicture}%
}

% --- Commande FlecheEllipse finale ---
% #1 = nom du nœud source (A)
% #2 = angle sortie du nœud source
% #3 = angle entrée du nœud cible
% #4 = décalage horizontal dx (en cm)
% #5 = décalage vertical dy (en cm)
% #6 = couleur ellipses et flèche
% #7 = contenu du nœud B (texte, formule, TikZ, etc.)
\newcommand{\FlecheEllipse}[7]{%
  \begin{tikzpicture}[overlay]
    % --- Nœud source A ---
    \node[draw,#6,ellipse,inner sep=2pt](#1){};

    % --- Nœud cible B ---
    \coordinate (B) at ($(#1) + (#4cm,#5cm)$);
    \node[draw,#6,ellipse,inner sep=2pt] (B-node) at (B) {#7};

    % --- Flèche de A vers B ---
    \draw[->,#6,thick] (#1) edge [out=#2,in=#3] (B-node);
  \end{tikzpicture}%
}

\renewcommand{\arraystretch}{0.2}

% Custom commands for events
\newcommand{\cvdegree}[6]{%
  \textbf{#1} & \textbf{#2} -- #3 \\
  & \textit{#4} {\color{headerblue}} \\
  & #5 \\[0.6em]
}

\newcommand{\cvevent}[6]{%
  \textbf{#1} & \textbf{#2} -- #3 \\
  & \textit{#4} {\color{cvred}} \\
  & #5 \\[0.6em]
}

\tcbuselibrary{skins}

\NewTotalTColorBox[auto counter]{\Example}{ +m +m }{
  enhanced,
  notitle,
  colback=white,
  colbacklower=white,
  frame hidden,
  boxrule=0pt,
  bicolor,
  sharp corners,
  %borderline west={4pt}{0pt}{greenForet},
  left=3mm, fontupper=\sffamily, overlay={% espace pour le texte vertical
      % Bande verticale
      \fill[greenForet]
      (frame.south west) rectangle
      ([xshift=4pt]frame.north west);

      % Texte vertical
      \node[
        rotate=90,
        anchor=south,
        font=\sffamily\bfseries,
        text=greenForet
      ] at ([xshift=2pt]frame.west) {#1};

      % Texte vertical centré
      \node[
        rotate=90,
        anchor=center,
        font=\sffamily\bfseries,
        text=white
      ] at ([xshift=2pt]frame.west |- frame.center) {Example};
    }
}{
  \textcolor{greenForet}{
    \sffamily\textbf{Exemple~: #1}\\
  }
  #2
}

\newcommand{\EncadreDeux}[4]{%
  \begin{tcolorbox}[
      enhanced,
      notitle,
      colback=white,
      colbacklower=white,
      frame hidden,
      boxrule=0pt,
      sharp corners,
      borderline west={4pt}{0pt}{#3},
      left=7mm,
      fontupper=\sffamily,
      overlay={%
          % Bande verticale
          \fill[#3]
          (frame.south west) rectangle
          ([xshift=4pt]frame.north west);

          \node[
            rotate=90,
            anchor=south,
            font=\sffamily\bfseries,
            text=#3
          ] at ([xshift=2pt]frame.west) {#2};

          % Texte vertical (titre)
          \node[
            rotate=90,
            anchor=center,
            font=\sffamily\bfseries\small,
            text=white
          ] at ([xshift=2pt]frame.west |- frame.center)
          {#1};

        }
    ]
    % contenu
    %\textcolor{#3}{\sffamily\textbf{#1}}\\
    #4
  \end{tcolorbox}
}

\NewTotalTColorBox[auto counter]{\Information}{+m}{
  enhanced,
  notitle,
  colback=white,
  colbacklower=white,
  frame hidden,
  boxrule=0pt,
  sharp corners,
  left=3mm,
  fontupper=\sffamily,
  overlay={
      % Bande verticale
      \fill[grayPierre]
      (frame.south west) rectangle
      ([xshift=4pt]frame.north west);

      % Texte vertical
      % \node[
      %     rotate=90,
      %     anchor=south,
      %     font=\sffamily\bfseries,
      %     text=grayPierre
      % ] at ([xshift=2pt]frame.west) {Information};

      % Texte vertical centré
      \node[
        rotate=90,
        anchor=center,
        font=\sffamily\bfseries,
        text=white
      ] at ([xshift=2pt]frame.west |- frame.center) {Information};
    },
}{
  #1
}

\NewTotalTColorBox[auto counter]{\ARetenir}{+m}{
  enhanced,
  notitle,
  colback=white,
  colbacklower=white,
  frame hidden,
  boxrule=0pt,
  bicolor,
  sharp corners,
  %borderline west={4pt}{0pt}{woodChene},
  left=3mm, fontupper=\sffamily, overlay={
      % Bande verticale
      \fill[woodChene]
      (frame.south west) rectangle
      ([xshift=4pt]frame.north west);

      % Texte vertical
      % \node[
      %     rotate=90,
      %     anchor=south,
      %     font=\sffamily\bfseries,
      %     text=woodChene
      % ] at ([xshift=2pt]frame.west) {À Retenir};

      % Texte vertical centré
      \node[
        rotate=90,
        anchor=center,
        font=\sffamily\bfseries,
        text=white
      ] at ([xshift=2pt]frame.west |- frame.center) {À Retenir};
    },
}{
  % \textcolor{woodChene}{
  %     \sffamily\textbf{À retenir~: #1}
  % }
  #1
}

% \newcommand{\EncadreUn}[3]{

% \NewTotalTColorBox[auto counter]{\ARetenir}{+m}{
%   enhanced,
%   notitle,
%   colback=white,
%   colbacklower=white,
%   frame hidden,
%   boxrule=0pt,
%   bicolor,
%   sharp corners,
%   borderline west={4pt}{0pt}{#3},
%   left=3mm, fontupper=\sffamily, overlay={
%       % Bande verticale
%       \fill[#3]
%       (frame.south west) rectangle
%       ([xshift=4pt]frame.north west);

%       % Texte vertical centré
%       \node[
%         rotate=90,
%         anchor=center,
%         font=\sffamily\bfseries,
%         text=white
%       ] at ([xshift=2pt]frame.west |- frame.center) {#1};
%     },
% }{
%   #2
% }

% }

\newcommand{\EncadreUn}[3]{%
  \begin{tcolorbox}[
      enhanced,
      notitle,
      colback=white,
      colbacklower=white,
      frame hidden,
      boxrule=0pt,
      sharp corners,
      borderline west={4pt}{0pt}{#2},
      left=6mm,
      fontupper=\sffamily,
      overlay={
          % Bande verticale
          \fill[#2]
          (frame.south west) rectangle
          ([xshift=4pt]frame.north west);

          % Texte vertical
          \node[
            rotate=90,
            anchor=center,
            font=\sffamily\bfseries\small,
            text=white
          ] at ([xshift=2pt]frame.west |- frame.center)
          {#1};
        }
    ]
    #3
  \end{tcolorbox}
}

\NewTotalTColorBox[auto counter]{\Rappel}{+m}{
  enhanced,
  notitle,
  colback=white,
  colbacklower=white,
  frame hidden,
  boxrule=0pt,
  bicolor,
  sharp corners,
  %borderline west={4pt}{0pt}{woodNoyer},
  left=3mm, fontupper=\sffamily, overlay={
      % Bande verticale
      \fill[woodNoyer]
      (frame.south west) rectangle
      ([xshift=4pt]frame.north west);

      % Texte vertical
      % \node[
      %     rotate=90,
      %     anchor=south,
      %     font=\sffamily\bfseries,
      %     text=woodNoyer
      % ] at ([xshift=2pt]frame.west) {Rappel};

      % Texte vertical centré
      \node[
        rotate=90,
        anchor=center,
        font=\sffamily\bfseries,
        text=white
      ] at ([xshift=2pt]frame.west |- frame.center) {Rappel};
    },
}{
  % \textcolor{woodNoyer}{
  %     \sffamily\textbf{Rappel~: #1}
  % }
  #1
}

\NewTotalTColorBox[auto counter]{\Resume}{+m}{
  enhanced,
  notitle,
  colback=white,
  colbacklower=white,
  frame hidden,
  boxrule=0pt,
  bicolor,
  sharp corners,
  %borderline west={4pt}{0pt}{woodAcajou},
  left=3mm, fontupper=\sffamily, overlay={
      % Bande verticale
      \fill[woodAcajou]
      (frame.south west) rectangle
      ([xshift=4pt]frame.north west);

      % Texte vertical
      % \node[
      %     rotate=90,
      %     anchor=south,
      %     font=\sffamily\bfseries,
      %     text=woodAcajou
      % ] at ([xshift=2pt]frame.west) {Conclusion};

      % Texte vertical centré
      \node[
        rotate=90,
        anchor=center,
        font=\sffamily\bfseries,
        text=white
      ] at ([xshift=2pt]frame.west |- frame.center) {Conclusion};
    },
}{
  % \textcolor{woodAcajou}{
  %     \sffamily\textbf{Conclusion~: #1}
  % }
  #1
}

% Définition de la fonction pour générer la grille
% \newcommand{\drawgrid}[4]{

%   \begin{scope}[opacity = 0.2]

%     % Paramètres : #1=xmin, #2=xmax, #3=ymin, #4=ymax

%     % Grille principale en cm
%     \draw[very thin, gray] (#1,#3) grid (#2,#4); % Grille de base (1 cm)

%     \draw[very thin, gray] (#1,#3) grid[step=0.1] (#2,#4);

%     \draw[thin, black] (#1,#3) edge[thick,line width=0.8ex,->,>=Stealth ] (#2,#3);

%     \foreach \x in {#1,..., #2} {
%         \draw[shift={(\x, #3)}] (0,-0.1) edge[line width=0.8ex] node[pos = -1 , below] {\x} (0,0.1); % Lignes verticales en mm
%       }

%     \draw[thin, black] (#1,#3)edge[thick,line width=0.8ex,->,>=Stealth ] (#1,#4);

%     \foreach \y in {#3,..., #4} {
%         \draw[shift={(#1, \y)}] (-0.1 , 0 ) edge[line width=0.8ex] node[pos = -1 , left ] {\y} (0.1 , 0); % Lignes verticales en mm
%       }

%   \end{scope}

% }

% Définition de la fonction pour générer la grille
\newcommand{\drawgrid}[4]{

  % #1 xmin
  % #2 xmax
  % #3 ymin
  % #4 ymax

  \begin{scope}[opacity=0.2]

    % Taille des axes
    \pgfmathsetmacro{\dx}{#2-#1}
    \pgfmathsetmacro{\dy}{#4-#3}

    % ======================
    % Grille principale
    % ======================

    \draw[very thin, gray]
    (#1,#3) grid (#2,#4);

    \draw[very thin, gray]
    (#1,#3) grid[step=0.1] (#2,#4);

    % ======================
    % Grille fine conditionnelle
    % ======================

    \ifdim \dx pt < 0.5pt

      \draw[very thin, gray , line width=0.01mm ]
      (#1,#3) grid[step=0.01] (#2,#4);

    \fi

    \ifdim \dy pt < 0.5pt

      \draw[very thin, gray , line width=0.01mm]
      (#1,#3) grid[step=0.01] (#2,#4);

    \fi

    % ======================
    % Axe x
    % ======================

    \draw[thick,line width=0.8ex,->,>=Stealth]
    (#1,#3) -- (#2,#3);

    \foreach \x in {#1,...,#2} {

        \draw[shift={(\x,#3)}]
        (0,-0.1) --
        (0,0.1)

        node[pos=-1,below] {\x};

      }

    % ======================
    % Axe y
    % ======================

    \draw[thick,line width=0.8ex,->,>=Stealth]
    (#1,#3) -- (#1,#4);

    \foreach \y in {#3,...,#4} {

        \draw[shift={(#1,\y)}]
        (-0.1,0) --
        (0.1,0)

        node[pos=-1,left] {\y};

      }

  \end{scope}

}